\section{Controls and robustness}

The controls change only one thing: the exact string scored after \texttt{"Final answer: "}. Figure~\ref{fig:controls} shows why that matters. Under identity ordering, the primary exact-option readout yields a \SemanticIdentityValidRate{} valid-rate for $\Delta_{\mathrm{margin}}$. If we keep the prompt fixed but score the literal continuation \texttt{"The answer is \{option\}"} instead, that valid-rate falls to \TemplatedIdentityValidRate{}. If we score only the answer letter, it rises to \LetterIdentityValidRate{}.

\begin{figure}[t]
\centering
\includegraphics[width=\columnwidth]{figure3_control_readout_comparison.pdf}
\caption{\textbf{Changing only the answer format changes the measurement.} Panel A shows how often $\Delta_{\mathrm{margin}}$ is defined for each readout. Panel B shows how often the two control readouts pick the same final winner as the semantic-exact primary readout on fully paired cases. Identity and reverse order behave similarly, which is good. The answer-format controls do not, which is exactly why they must stay controls rather than replacing the main semantic path.}
\label{fig:controls}
\end{figure}

The readouts also disagree on which answer the model prefers at the end. In fully paired identity cases, the templated control agrees with the exact-option path on only \TemplatedAgreementIdentity{} of cases, the letter-label control agrees on only \LetterAgreementIdentity{}, and all three readouts agree on just \AllSameReadoutRate{} of cases. This does not weaken the main result. It tells us something important about measurement: scoring \texttt{"A"} is not the same experiment as scoring \texttt{"photosynthesis"}. The controls are useful because they reveal answer-format effects. They are not suitable replacements for the main semantic readout.

The answer-order control tells a different story. In Figure~\ref{fig:controls}, the identity and reverse traces are nearly on top of each other. That is reassuring because it suggests the main effect is not coming from one arbitrary option order. The report-only robustness summaries also line up with the primary rule. The secondary criterion yields the same positive-versus-insufficient pattern as the primary rule (\OverallSecondaryPositiveCells{} positive, \OverallSecondaryInsufficientCells{} insufficient), and the gold-answer sensitivity analysis is positive on all \GoldArcCsqaPositiveCells{} evaluable ARC-Challenge and CommonsenseQA cells while remaining underpowered on the MMLU slice (\GoldMmluInsufficientCells{} insufficient cells).

Finally, the QC path is designed to surface problems, not hide them. The released audit bundle logs \TokenCollisionCount{} tokenization collisions and \DuplicateOptionCount{} duplicate-option examples instead of silently dropping them, while the merged run itself still has zero failure rows and zero unpaired rows. The point is not to make the result look tidy. The point is to show that the result still holds when we keep the pipeline honest about edge cases.

\FloatBarrier
